% =====================================================================
%  ANTES — BROCHURA COM ORELHAS — variante 100% PRETO E BRANCO
%  Idêntico a brochura-template.tex, exceto pelas três cores nomeadas.
%  Engine: XeLaTeX
%  Body  : TeX Gyre Pagella 11/14
%  Trim  : 16 × 23 cm (160 × 230 mm) + 2,5 mm de sangria → página 165 × 235 mm
%  Color : miolo inteiro em cinza — tipografia E figuras. A cor vive só na capa.
%  Sans  : Inter (titles, captions)
%  Lang  : Brazilian Portuguese (polyglossia)
% =====================================================================
\documentclass[11pt,openright,twoside]{book}

% --- Página = trim 16×23 cm + 2,5 mm de sangria em cada lado = 165×235 mm ---
% As margens são medidas a partir da borda da PÁGINA, então já incluem a faixa
% de sangria. A menor margem usada aqui é 22 mm (= 2,5 de sangria + 19,5 de
% recuo), o que mantém ≥ 0,5 cm de segurança da linha de corte.
% bindingoffset maior que o da capa dura: lombada colada (hot melt) abre menos
% que costurada, então o miolo precisa de mais medianiz.
\usepackage[
  paperwidth=165mm, paperheight=235mm,
  inner=22mm,
  outer=18mm,
  top=22mm,
  bottom=22mm,
  headsep=5mm,
  footskip=8mm,
  bindingoffset=6mm,
]{geometry}

% --- Language / typography ---
\usepackage{polyglossia}
\setdefaultlanguage[variant=brazilian]{portuguese}
\usepackage{fontspec}
\usepackage{microtype}

% Body: TeX Gyre Pagella (Palatino clone) — full Latin coverage, classic book look.
% EB Garamond on Debian has incomplete Bold variant for accented Portuguese.
\setmainfont{TeX Gyre Pagella}[
  Ligatures={Common,TeX}
]
\setsansfont{Inter}[
  Scale=0.92,
  Numbers={Lining,Proportional}
]
\newfontfamily\dropfont{TeX Gyre Pagella}

% --- Glifos que a Pagella não tem -------------------------------------------
% TeX Gyre Pagella não traz U+2080..U+2089 (dígitos subscritos). Sem
% mapeamento o XeLaTeX imprime o .notdef box e "VO₂ max" sai como "VO□ max"
% — foram 31 ocorrências na prova de 2026-08-24. Mapear para \textsubscript
% compõe o dígito na PRÓPRIA Pagella, rebaixado e reduzido: fica correto
% tipograficamente e não puxa fonte de fallback (que mudaria a cor do texto).
% Superscritos idem, por simetria — a fonte tem ¹²³ mas não ⁴..⁹.
\usepackage{newunicodechar}
\newunicodechar{₀}{\textsubscript{0}}
\newunicodechar{₁}{\textsubscript{1}}
\newunicodechar{₂}{\textsubscript{2}}
\newunicodechar{₃}{\textsubscript{3}}
\newunicodechar{₄}{\textsubscript{4}}
\newunicodechar{₅}{\textsubscript{5}}
\newunicodechar{₆}{\textsubscript{6}}
\newunicodechar{₇}{\textsubscript{7}}
\newunicodechar{₈}{\textsubscript{8}}
\newunicodechar{₉}{\textsubscript{9}}
\newunicodechar{⁴}{\textsuperscript{4}}
\newunicodechar{⁵}{\textsuperscript{5}}
\newunicodechar{⁶}{\textsuperscript{6}}
\newunicodechar{⁷}{\textsuperscript{7}}
\newunicodechar{⁸}{\textsuperscript{8}}
\newunicodechar{⁹}{\textsuperscript{9}}

% Títulos e legendas são ragged-right: hifenizar ali não economiza nada, só
% deixa a palavra partida no meio de um display. \notitlehyphen desliga a
% hifenização no escopo do título (não vaza para o corpo).
\newcommand{\notitlehyphen}{\hyphenpenalty=10000\exhyphenpenalty=10000\relax}

% --- Colors (Plenya palette) ---
\usepackage[table]{xcolor}   % [table] habilita \rowcolor (tabelas nativas)
% Variante P&B: as três cores viram cinza. Não é dessaturação por luminância —
% petrol (#063b4f) tem luminância 49, o que daria #313131, MAIS CLARO que o
% corpo de texto (#1a1a1a), deixando o título mais fraco que o parágrafo. No
% impresso a hierarquia tem de vir do peso e do corpo da fonte, então petrol
% vira quase preto e o gold vira um cinza claramente acessório.
\definecolor{petrol}{HTML}{111111}   % títulos, capitulares, rótulo de legenda
\definecolor{gold}{HTML}{767676}     % numeração "PARTE N" / "CAPÍTULO N"
\definecolor{ink}{HTML}{1a1a1a}      % corpo — inalterado

% --- Drop caps ---
\usepackage{lettrine}
\renewcommand{\LettrineFontHook}{\dropfont\color{petrol}}
\setlength{\DefaultNindent}{0pt}
\renewcommand{\DefaultLraise}{0.05}

% --- Title structure ---
\usepackage{titlesec}
\titleformat{\part}[display]
  {\sffamily\filcenter\color{petrol}\notitlehyphen}
  {\normalfont\sffamily\color{gold}\fontsize{12}{14}\selectfont\MakeUppercase{Parte~\thepart}}
  {2em}
  {\fontsize{38}{44}\selectfont\bfseries}
\titlespacing*{\part}{0pt}{80pt}{40pt}

\titleformat{\chapter}[display]
  {\sffamily\bfseries\color{petrol}\filright\notitlehyphen}
  {\normalfont\sffamily\color{gold}\fontsize{12}{14}\selectfont CAPÍTULO~\thechapter}
  {1.2em}
  {\fontsize{28}{34}\selectfont}
\titlespacing*{\chapter}{0pt}{30pt}{30pt}

\titleformat{\section}
  {\sffamily\bfseries\color{petrol}\large\notitlehyphen}{\thesection}{0.6em}{}
\titleformat{\subsection}
  {\sffamily\bfseries\color{ink}\normalsize\notitlehyphen}{\thesubsection}{0.6em}{}

% \titlespacing SEM estrela = o primeiro parágrafo depois do título entra com
% recuo. É a norma brasileira de diagramação e o editor apontou a ausência em
% quatro trechos diferentes (2026-08-26). No \chapter continua COM estrela: ali
% quem abre é a capitular, que não leva recuo.
\titlespacing{\section}{0pt}{3.5ex plus 1ex minus .2ex}{2.3ex plus .2ex}
\titlespacing{\subsection}{0pt}{3.25ex plus 1ex minus .2ex}{1.5ex plus .2ex}

% --- Headers / footers ---
\usepackage{fancyhdr}
\setlength{\headheight}{14pt}
\pagestyle{fancy}
\fancyhf{}
\renewcommand{\headrulewidth}{0pt}
% Recto and verso both show the chapter title (no section names).
% \leftmark holds the chapter title (set by \chaptermark below); \rightmark
% would normally hold the section title — we override \sectionmark to a no-op
% so the chapter title persists across the whole chapter on every page.
\fancyhead[LE]{\sffamily\footnotesize\color{ink}\thepage\quad\textcolor{gray}{·}\quad ANTES}
\fancyhead[RO]{\sffamily\footnotesize\color{ink}\nouppercase{\leftmark}\quad\textcolor{gray}{·}\quad\thepage}
% Chapter mark for headers: keep only the short title (text before the
% first em-dash), to avoid long chapter titles overflowing the right margin.
\makeatletter
\def\@firstdash#1 — #2\@nil{#1}
\renewcommand{\chaptermark}[1]{\markboth{\@firstdash#1 — \@nil}{}}
\makeatother
\renewcommand{\sectionmark}[1]{}

\fancypagestyle{plain}{\fancyhf{}\renewcommand{\headrulewidth}{0pt}\renewcommand{\footrulewidth}{0pt}}

% --- Spacing / paragraphs ---
% Medida de linha de 125 mm (4,92") pede corpo 11pt: 10,5pt daria ~75
% caracteres por linha, acima da faixa confortável de leitura.
\linespread{1.10}
\setlength{\parskip}{0pt}
\setlength{\parindent}{1.2em}
% ATENÇÃO: este \frenchspacing NÃO pega. O polyglossia troca o idioma dentro
% de \begin{document} e restaura o \nonfrenchspacing (sfcode 3000 depois de
% ponto, 2000 depois de dois-pontos), então o miolo sai com espaço dobrado no
% fim de frase — medido em 2026-08-27: 4,33pt depois de ponto contra 3,27pt
% entre palavras. Um \frenchspacing logo depois do \begin{document} corrige,
% mas mexe no cinza de TODAS as páginas e custa hífens (o espaço extra é
% elasticidade que o TeX usa para fechar a linha sem hifenizar: 339 → 424).
% Fica como está até o autor decidir — e, se decidir, vale para as três
% edições juntas (KDP, capa dura e brochura), não só para esta.
\frenchspacing
% Penalidades de quebra. Segunda rodada, medida em 2026-08-27 sobre a prova
% da brochura com build/analyze-hyphenation.py (varredura BASE/A..E; a rodada
% de 2026-08-23 tinha deixado 8,2% das linhas do corpo terminando em hífen,
% que o autor recusou). Os números abaixo são do miolo P&B, 352 páginas:
%
%              hífens no corpo   espaço p99   linhas >2,5x   páginas
%   antes           8,2%            2,05          0,1%         352
%   agora           3,5%            2,46          0,5%         352
%
%   brokenpenalty=10000  proíbe hífen na ÚLTIMA linha da página — a sílaba
%     ficaria órfã na página seguinte.
%   hyphenpenalty=8000 (era 1000) é o que derruba a taxa pela metade. Acima
%     disso (10000 = hifenização desligada) a conta vira ruim: 0% de hífens,
%     mas 22% das linhas frouxas e 4% delas acima de 2,5x o espaço natural —
%     rio branco visível na página. 8000 mantém a cauda em 0,5%.
%   exhyphenpenalty=10000 proíbe quebrar DEPOIS de um hífen que já existe na
%     palavra. É o que tirava "check-/up", "pós-/graduação", "terça-/feira":
%     em português composto isso é o hífen mais feio dos que sobravam.
%   lefthyphenmin=3 / righthyphenmin=4 (era 2/3) removem os pedaços curtos
%     ("in-", "as-", "-me") — os cortes que mais custam ao leitor.
%   emergencystretch=1.5em: com 1em o TeX não completa as linhas que sobram
%     sem hifenizar; com 2em ele ganha licença demais e a cauda estoura.
\widowpenalty=10000
\clubpenalty=10000
\brokenpenalty=10000
\hyphenpenalty=8000
\exhyphenpenalty=10000
\doublehyphendemerits=10000
\finalhyphendemerits=10000
% \lefthyphenmin / \righthyphenmin vivem depois do \begin{document} (ver lá):
% o polyglossia troca o idioma dentro do \begin{document} e restaura os valores
% do padrão de hifenização (2/3), apagando qualquer ajuste feito aqui no
% preâmbulo. Medido em 2026-08-27 com \typeout.
\tolerance=3000
\emergencystretch=1.5em

% --- Quotes / blockquote ---
\usepackage{quoting}
\quotingsetup{
  vskip=0.6\baselineskip,
  leftmargin=1.2em,
  rightmargin=0pt,
  font=itshape,
}
\renewenvironment{quote}{\begin{quoting}}{\end{quoting}}

% --- Tables (calc required by Pandoc 3.x p{} column-spec arithmetic) ---
\usepackage{calc}
\usepackage{longtable}
\usepackage{booktabs}
\usepackage{array}
\renewcommand{\arraystretch}{1.15}

% --- Images ---
\usepackage{graphicx}
\usepackage{float}
% Com os defaults do book (\topfraction 0.7, \textfraction 0.2,
% \floatpagefraction 0.5) uma figura de meia página não cabia junto do texto e
% ia parar numa página de float só dela, com dois terços em branco — e caía
% duas páginas depois do parágrafo que a chama. O editor apontou os dois
% efeitos (2026-08-26: "ficou totalmente isolado porque se refere ao texto da
% página 36"). Afrouxar as frações deixa a figura ficar na página do texto.
\renewcommand{\topfraction}{0.92}
\renewcommand{\bottomfraction}{0.85}
\renewcommand{\textfraction}{0.08}
\renewcommand{\floatpagefraction}{0.80}
\setcounter{topnumber}{3}
\setcounter{bottomnumber}{2}
\setcounter{totalnumber}{4}
\floatplacement{figure}{htbp}
% Cap every \includegraphics at the printable area so figures never bleed past
% the page. keepaspectratio preserves geometry; width=\linewidth shrinks wide
% figures; height limit prevents tall portrait figures from overflowing.
\setkeys{Gin}{width=\linewidth, height=0.82\textheight, keepaspectratio}
\usepackage{caption}
\DeclareCaptionFont{nohyph}{\notitlehyphen}
\captionsetup{
  font={small,sf,nohyph},
  labelfont={bf,color=petrol},
  labelsep=period,
  justification=justified,   % editor 2026-08-26: legenda de figura
                             % alinhada à esquerda destoa do miolo
  singlelinecheck=false,
  skip=4pt
}

% --- Hyperref last ---
\usepackage[unicode]{hyperref}
\hypersetup{
  colorlinks=false,
  pdfborder={0 0 0},
  pdftitle={ANTES — A Janela Silenciosa entre o Normal e o Ótimo — onde a saúde é decidida},
  pdfauthor={Dr. Getulio Amaral Filho},
  pdfsubject={Medicina preventiva e longevidade — brochura com orelhas, miolo P\&B},
  pdfkeywords={longevidade, biomarcadores, AGIR, medicina preventiva, brochura}
}

% --- Pandoc compatibility ---
% \textsubscript/\textsuperscript vindos do \newunicodechar acima não podem
% entrar cru na string de bookmark do PDF — hyperref reclama. Na outline o
% dígito vira texto normal ("VO2 Max"), que é o comportamento certo.
\pdfstringdefDisableCommands{%
  \def\textsubscript#1{#1}%
  \def\textsuperscript#1{#1}%
}

\providecommand{\tightlist}{\setlength{\itemsep}{0pt}\setlength{\parskip}{0pt}}
\providecommand{\pandocbounded}[1]{#1}
\providecommand{\real}[1]{#1}

% --- TOC depth: chapters only (no sections / subsections) ---
% Per author's call (2026-04-25): keep TOC short, "Outlive-like" — readers folheiam
% para ver o arco do livro, não para navegar microestrutura. Sections still numbered
% in the body via secnumdepth=2 below.
\setcounter{tocdepth}{0}      % \chapter (and \part) only in TOC
\setcounter{secnumdepth}{2}   % keep numbering for sections + subsections in body

% --- TOC formatting: long chapter titles wrap to next line and respect
%     minimum spacing before the page number. Reserves 2.5em on the right
%     for the page number ("rmargin"), so even a 3-digit page (e.g. 411)
%     never collides with the title text.
\usepackage{titletoc}
\titlecontents{chapter}[2.4em]                          % left indent for title body
  {\addvspace{0.7em}\bfseries\sffamily\color{petrol}\notitlehyphen\raggedright}
  {\contentslabel[\thecontentslabel]{2.4em}}             % numbered label box
  {\hspace*{-2.4em}}                                     % unnumbered (Introdução etc.)
  % \raggedright acima deixa a 1ª linha de um título de duas linhas em bandeira
  % em vez de esticada de margem a margem. O preenchimento até o número da
  % página tem de ser de ordem MAIOR (filll) que a elasticidade do \rightskip
  % (fil) do \raggedright, senão os dois empatam e o número desgruda da margem.
  {\hspace{1em}\nobreak\hskip 0pt plus 1filll\nobreak\contentspage}
  [\vspace{0.1em}]                                       % space after

% (right margin reserved for page numbers is handled inside titlecontents above
% via \hspace{1em}\hfill — adequate breathing room without touching @tocrmarg)

% --- Pandoc title vars hook ---
$if(title)$
\title{$title$}
$endif$
$if(author)$
\author{$for(author)$$author$$sep$ \and $endfor$}
$endif$

\begin{document}
% Tem de ser AQUI, depois que o polyglossia já escolheu o idioma. 3/4 remove os
% cortes curtos ("in-", "as-", "-me"), os que mais custam ao leitor.
\lefthyphenmin=3
\righthyphenmin=4

% =====================================================================
%  FRONT MATTER (Roman numerals i, ii, iii…)
% =====================================================================
\frontmatter
\pagestyle{empty}

% --- HALF TITLE (i) ---
\thispagestyle{empty}
\vspace*{4cm}
\begin{center}
{\sffamily\fontsize{36}{42}\selectfont\color{petrol} ANTES}
\end{center}
\clearpage

% --- BLANK (ii) ---
\thispagestyle{empty}\null\clearpage

% --- TITLE PAGE (iii) ---
\thispagestyle{empty}
\vspace*{3cm}
\begin{center}
{\sffamily\fontsize{52}{58}\selectfont\color{petrol} ANTES}\\[2.2em]
{\itshape\fontsize{16}{22}\selectfont A Janela Silenciosa\\[0.1em]
entre o Normal e o Ótimo\\[0.3em]
— onde a saúde é decidida}\\[5em]
{\sffamily\fontsize{14}{18}\selectfont\color{ink} Dr. Getulio Amaral Filho}\\[6em]
\vfill
{\sffamily\footnotesize\color{gold} PLENYA}\\[0.4em]
{\sffamily\footnotesize\color{ink} Edição do Autor \quad·\quad 2026}
\end{center}
\clearpage

% --- Body of frontmatter (copyright, dedication, epigraph, TOC) comes from Pandoc ---
$body$

\end{document}
